\documentclass[sigconf, nonacm]{acmart}

% ============================================================================
% PRISMA SUMMER SCHOOL 2026 -- GROUP WORK REPORT
% ============================================================================
\setcopyright{none}
\settopmatter{printacmref=false, printccs=false, printfolios=true}
\acmDOI{}
\acmISBN{}
\acmConference[PRISMA '26]{PRISMA Summer School on Uncertainty in Energy Systems}{June 15--19, 2026}{University of Geneva, Switzerland}

% ============================================================================
% PACKAGES
% ============================================================================
\usepackage{amsmath,amsfonts}
\usepackage{booktabs}
\usepackage{siunitx}
\usepackage{graphicx}
\usepackage{float}

\sisetup{per-mode=symbol, exponent-product=\times, group-separator={,}}

% ============================================================================
\begin{document}
% ============================================================================

\title{Probabilistic Methods Applied to the Scenario Compass Initiative: What Does It Take to Reach Net Zero CO$_2$ by 2070?}

\author{Lucas Alvarez}
\email{lualpi@upv.es}
\affiliation{%
  \institution{Universitat Polit\`ecnica de Val\`encia}
  \city{Valencia}
  \country{Spain}
}

\author{Matteo Catania}
\email{matteo.catania@polimi.it}
\affiliation{%
  \institution{Politecnico di Milano}
  \city{Milan}
  \country{Italy}
}

\author{R\'emi Paccou}
\email{remi.paccou@proton.me}
\affiliation{%
  \institution{CIRED}
  \city{Nogent-sur-Marne}
  \country{France}
}
\affiliation{%
  \institution{Schneider Electric}
  \city{Grenoble}
  \country{France}
}

% ----------------------------------------------------------------------------
\begin{abstract}
% TODO
\end{abstract}

\keywords{probabilistic methods, scenario ensemble, net zero, Scenario Compass Initiative, uncertainty}

\maketitle

% ============================================================================
\section{Introduction}\label{sec:intro}
% ============================================================================

The Scenario Compass Initiative (SCI) assembles 1,564 pathways from 65 integrated assessment models, of which 497 reach net zero CO$_2$ emissions globally by 2070. This ensemble encodes substantial structural disagreement across models regarding technology costs, deployment rates, and policy assumptions. The question we address is: \emph{which variable, when modeled probabilistically, best discriminates pathways that achieve net zero by 2070 from those that do not?}

Our approach proceeds in four steps. First, we partition the ensemble into two groups based on the net-zero-by-2070 criterion. Second, we compare the cross-scenario distributions of candidate variables between the two groups, using effect-size measures to rank discriminatory power. Third, we calibrate a simple probabilistic model on the selected variable. Fourth, we derive a conditional probability of reaching net zero as a function of that variable's trajectory.

% ============================================================================
\section{Data and Methods}\label{sec:methods}
% ============================================================================

% TODO

% ============================================================================
\section{Results}\label{sec:results}
% ============================================================================

% TODO

% ============================================================================
\section{Discussion and Conclusions}\label{sec:discussion}
% ============================================================================

% TODO

% ============================================================================
\bibliographystyle{ACM-Reference-Format}
\begin{thebibliography}{10}

% TODO

\end{thebibliography}

\end{document}
